% GENERATED FILE. Edit canonical lessons, then run npm run generate:books.
% canonical-source-hash: fnv1a64:3110e780b00030e0
% canonical-lessons: ZH-C396-buguan, ZH-C396-jiran, ZH-C396-chufei, ZH-C396-quantian, ZH-C396-jishi, ZH-R396-review-392-396

\chapter{Whatever Happens}
\label{ch:zh-whatever-happens}
\hlchaptermodality{396}

\begin{chapteropening}
\textbf{By the end of this chapter:} \emph{I can talk about doing something no matter what, since something is so, unless something happens, all day and even if something happens.}

Complete the last lesson of chapter 396: I can talk about doing something no matter what, since something is so, unless something happens, all day and even if something happens.
\end{chapteropening}

\section[bùguǎn]{bùguǎn — no matter (whether or how)}

\label{lesson:ZH-C396-buguan}



\begin{quote}
Before the new one: say the Mandarin for to get (a job) done, then the Mandarin for a condition (that must be met).
\end{quote}

\subsection*{You'll want to know: bùguǎn}
\emph{bùguǎn} — \textquotedblleft{}no matter (whether or how)\textquotedblright{}.

\begin{quote}
\emph{bùguǎn míngtiān tiānqì zěnmeyàng wǒ dōu qù} — Whatever the weather tomorrow, I'm going.
\end{quote}

Its characters come later; for now it is a word for the ear.

Tone: \zh{不} keeps its own falling tone here, bù, because guǎn after it is a third tone, not a falling one.

\begin{cousinweb}
\emph{bùguǎn} is built from \zh{不} 'not' and \emph{guǎn} 'to be in charge', both inherited from Middle Chinese, \emph{pjut} and \emph{kwanX} (Guangyun, in Baxter's transcription), which glosses the second 'a pipe' and 'to take charge': not caring which.
\end{cousinweb}

\begin{grammarlens}[title={What you've built}]
One more word for saying what would happen.
\end{grammarlens}

\subsection*{Your turn}
\begin{itemize}
\raggedright
  \item \emph{Say it:} \emph{bùguǎn}
  \item \emph{Say it:} \emph{bùguǎn}, once more
  \item \emph{Say it:} \emph{bùguǎn míngtiān tiānqì zěnmeyàng wǒ dōu qù} — Whatever the weather tomorrow, I'm going.
  \item \emph{Recall:} say the Mandarin for to get (a job) done, then the Mandarin for a condition (that must be met), then say \emph{bùguǎn} again
\end{itemize}

\begin{tcolorbox}[breakable,colback=teal!4,colframe=teal!35!black,title={Before you move on}]
What does bùguǎn mean? (\textquotedblleft{}No matter (whether or how)\textquotedblright{}.) Say it once more.
\end{tcolorbox}
\section[jìrán]{jìrán — given that; now that}

\label{lesson:ZH-C396-jiran}



\begin{quote}
Before the new one: say the Mandarin for a condition (that must be met), then the Mandarin for no matter (whether or how).
\end{quote}

\subsection*{You'll want to know: jìrán}
\emph{jìrán} — \textquotedblleft{}given that; now that\textquotedblright{}.

\begin{quote}
\emph{jìrán nǐ bìng le jiù zài jiā xiūxi} — Since you're ill, stay at home and rest.
\end{quote}

Its characters come later; for now it is a word for the ear.

\begin{cousinweb}
\emph{jìrán} is built from \emph{jì} 'already' and \zh{然} 'so', both inherited from Middle Chinese, \emph{kj+jH} and \emph{nyen} (Guangyun, in Baxter's transcription), which glosses the first 'finished' and the second 'like this'.
\end{cousinweb}

\begin{grammarlens}[title={What you've built}]
One more word for saying what would happen.
\end{grammarlens}

\subsection*{Your turn}
\begin{itemize}
\raggedright
  \item \emph{Say it:} \emph{jìrán}
  \item \emph{Say it:} \emph{jìrán}, once more
  \item \emph{Say it:} \emph{jìrán nǐ bìng le jiù zài jiā xiūxi} — Since you're ill, stay at home and rest.
  \item \emph{Recall:} say the Mandarin for a condition (that must be met), then the Mandarin for no matter (whether or how), then say \emph{jìrán} again
\end{itemize}

\begin{tcolorbox}[breakable,colback=teal!4,colframe=teal!35!black,title={Before you move on}]
What does jìrán mean? (\textquotedblleft{}Given that; now that\textquotedblright{}.) Say it once more.
\end{tcolorbox}
\section[chúfēi]{chúfēi — unless}

\label{lesson:ZH-C396-chufei}



\begin{quote}
Before the new one: say the Mandarin for no matter (whether or how), then the Mandarin for given that.
\end{quote}

\subsection*{You'll want to know: chúfēi}
\emph{chúfēi} — \textquotedblleft{}unless\textquotedblright{}.

\begin{quote}
\emph{chúfēi nǐ bìng le bùrán míngtiān yídìng yào lái} — Unless you're ill, be sure to come tomorrow.
\end{quote}

Its characters come later; for now it is a word for the ear.

\begin{cousinweb}
\emph{chúfēi} is built from \emph{chú} 'to take away' and \emph{fēi} 'not', both inherited from Middle Chinese, \emph{drjo} and \emph{pj+j} (Guangyun, in Baxter's transcription), which glosses the first 'to remove' and the second 'it is not so'.
\end{cousinweb}

\begin{grammarlens}[title={What you've built}]
One more word for saying what would happen.
\end{grammarlens}

\subsection*{Your turn}
\begin{itemize}
\raggedright
  \item \emph{Say it:} \emph{chúfēi}
  \item \emph{Say it:} \emph{chúfēi}, once more
  \item \emph{Say it:} \emph{chúfēi nǐ bìng le bùrán míngtiān yídìng yào lái} — Unless you're ill, be sure to come tomorrow.
  \item \emph{Recall:} say the Mandarin for no matter (whether or how), then the Mandarin for given that, then say \emph{chúfēi} again
\end{itemize}

\begin{tcolorbox}[breakable,colback=teal!4,colframe=teal!35!black,title={Before you move on}]
What does chúfēi mean? (\textquotedblleft{}Unless\textquotedblright{}.) Say it once more.
\end{tcolorbox}
\section[\texorpdfstring{\zh{全天} (quántiān)}{全天 (quántiān)}]{\zh{全天} (quántiān) — all day; round the clock}

\label{lesson:ZH-C396-quantian}



\begin{quote}
Before the new one: say the Mandarin for given that, then the Mandarin for unless.
\end{quote}

\subsection*{You'll want to know: \zh{全天}}
\textbf{\zh{全天}} — \emph{quántiān} — \textquotedblleft{}all day; round the clock\textquotedblright{}.

\begin{quote}
\textbf{\zh{要是晚上要用钱这儿全天都能取钱}} \emph{yàoshi wǎnshang yào yòng qián zhèr quántiān dōu néng qǔqián} — If you need cash at night, you can take money out here at any hour.
\end{quote}

\begin{cousinweb}
\textbf{\zh{全天}} (\emph{quántiān}) is built from \zh{全} 'whole' and \zh{天} 'a day', both inherited from Middle Chinese, \emph{dzjwen} and \emph{then} (Guangyun, in Baxter's transcription).
\end{cousinweb}

\begin{grammarlens}[title={What you've built}]
One more word for saying what would happen.
\end{grammarlens}

\subsection*{Your turn}
\begin{itemize}
\raggedright
  \item \emph{Say it:} \emph{quántiān}
  \item \emph{Say it:} \emph{quántiān}, once more
  \item \emph{Say it:} \emph{yàoshi wǎnshang yào yòng qián zhèr quántiān dōu néng qǔqián} — If you need cash at night, you can take money out here at any hour.
  \item \emph{Recall:} say the Mandarin for given that, then the Mandarin for unless, then say \emph{quántiān} again
\end{itemize}

\begin{tcolorbox}[breakable,colback=teal!4,colframe=teal!35!black,title={Before you move on}]
What does \zh{全天} mean? (\textquotedblleft{}All day; round the clock\textquotedblright{}.) Say it once more.
\end{tcolorbox}
\section[jíshǐ]{jíshǐ — even if}

\label{lesson:ZH-C396-jishi}



\begin{quote}
Before the new one: say the Mandarin for unless, then the Mandarin for round the clock.
\end{quote}

\subsection*{You'll want to know: jíshǐ}
\emph{jíshǐ} — \textquotedblleft{}even if\textquotedblright{}.

\begin{quote}
\emph{jíshǐ xiàyǔ wǒmen yě yào qù} — Even if it rains, we're still going.
\end{quote}

Its characters come later; for now it is a word for the ear.

\begin{cousinweb}
\emph{jíshǐ} is built from \emph{jí} 'to go up to' and \emph{shǐ} 'to make (someone do)', both inherited from Middle Chinese, \emph{tsik} and \emph{sriX} (Guangyun, in Baxter's transcription), which glosses the first 'to go to' and 'now' and the second 'to employ' and 'to order'. How the pair came to mean 'even if' is not settled.
\end{cousinweb}

\begin{grammarlens}[title={What you've built}]
One more word for saying what would happen.
\end{grammarlens}

\subsection*{Your turn}
\begin{itemize}
\raggedright
  \item \emph{Say it:} \emph{jíshǐ}
  \item \emph{Say it:} \emph{jíshǐ}, once more
  \item \emph{Say it:} \emph{jíshǐ xiàyǔ wǒmen yě yào qù} — Even if it rains, we're still going.
  \item \emph{Recall:} say the Mandarin for unless, then the Mandarin for round the clock, then say \emph{jíshǐ} again
\end{itemize}

\begin{tcolorbox}[breakable,colback=teal!4,colframe=teal!35!black,title={Before you move on}]
What does jíshǐ mean? (\textquotedblleft{}Even if\textquotedblright{}.) Say it once more.
\end{tcolorbox}
\section[(dialogue)]{Looking back: from A Change of Plan to Whatever Happens}

\label{lesson:ZH-R396-review-392-396}



\begin{quote}
Say the Mandarin for to alter and to change the date.
\end{quote}

\subsection*{Your turn}
\begin{itemize}
\raggedright
  \item \emph{Say it:} to alter, to change the date, to make different, to change the name (of), to change one's tune
  \item \emph{Say it:} to exhort (someone), be sure to (urging), all-purpose, at any time at all, going smoothly
  \item \emph{Say it:} to succeed, to fail, to take a risk, to be wasted (of effort), staying the same
  \item \emph{Say it:} to come true, to obtain (through effort), fruits of work, to get (a job) done, a condition (that must be met)
  \item \emph{Say it:} no matter (whether or how), given that, unless, round the clock, even if
  \item \emph{Read it:} \zh{改} \zh{万} \zh{取} — each one aloud, then a word that holds it
\end{itemize}

\begin{tcolorbox}[breakable,colback=teal!4,colframe=teal!35!black,title={Before you move on}]
To alter? (\textbf{\zh{改}}, \emph{gǎi}.) To change the date? (\textbf{\zh{改期}}, \emph{gǎiqī}.) To make different? (\textbf{\zh{改变}}, \emph{gǎibiàn}.) To change the name (of)? (\textbf{\zh{改名}}, \emph{gǎimíng}.) To change one's tune? (\textbf{\zh{改口}}, \emph{gǎikǒu}.) To exhort (someone)? (\emph{zhǔfu}.) Be sure to (urging)? (\textbf{\zh{千万}}, \emph{qiānwàn}.) All-purpose? (\textbf{\zh{万能}}, \emph{wànnéng}.) At any time at all? (\emph{suíshí}.) Going smoothly? (\emph{shùnlì}.) To succeed? (\emph{chénggōng}.) To fail? (\emph{shībài}.) To take a risk? (\emph{màoxiǎn}.) To be wasted (of effort)? (\emph{báifèi}.) Staying the same? (\textbf{\zh{不变}}, \emph{búbiàn}.) To come true? (\emph{shíxiàn}.) To obtain (through effort)? (\textbf{\zh{取得}}, \emph{qǔdé}.) Fruits of work? (\textbf{\zh{成果}}, \emph{chéngguǒ}.) To get (a job) done? (\emph{wánchéng}.) A condition (that must be met)? (\emph{tiáojiàn}.) No matter (whether or how)? (\emph{bùguǎn}.) Given that? (\emph{jìrán}.) Unless? (\emph{chúfēi}.) Round the clock? (\textbf{\zh{全天}}, \emph{quántiān}.) Even if? (\emph{jíshǐ}.)
\end{tcolorbox}
